\section{Further Analysis}
\subsection{Analysis of Historical Inductive Biases in Long-Horizon Discovery}
\definecolor{oursblue}{HTML}{0072B2}
\definecolor{baselinegray}{HTML}{6B7280}
\definecolor{guidanceorange}{HTML}{D55E00}
\definecolor{guidancegreen}{HTML}{009E73}

\begin{wrapfigure}[14]{r}{0.5\linewidth}
% \vspace{-1em}
\centering
\begin{tikzpicture}
\begin{axis}[
  width=\linewidth,
  height=4cm,
  xmin=0,
  xmax=1050,
  ymin=0.40,
  ymax=2.00,
  axis lines=left,
  axis line style={black!55,line width=0.65pt},
  xmajorgrids=false,
  ymajorgrids=true,
  grid style={black!8,line width=0.30pt},
  scaled ticks=false,
  tick align=outside,
  tick style={black!55,line width=0.50pt},
  tick label style={font=\scriptsize,text=black!80},
  xlabel style={font=\small,yshift=2pt},
  ylabel style={font=\small},
  clip=true,
  clip marker paths=true,
  xtick={0,250,500,750,1000},
  ytick={0.4,0.8,1.2,1.6,2.0},
  xlabel={Number of Generations},
  ylabel={Performance ($1/\mathrm{ms}$)},
  legend style={
    draw=none,
    font=\scriptsize,
    at={(0.43,1.03)},
    anchor=south,
    legend columns=2,
    row sep=1pt,
    column sep=0.8em
  }
]

% ============================================================
% Dream-RSI
% ============================================================
\addplot[
  const plot,
  color=oursblue,
  line width=2.0pt,
  mark=*,
  mark size=1.2pt,
  mark options={
    fill=oursblue,
    draw=oursblue
  }
]
table[x=budget,y=score]
{Main_Text/raw_data/data/convdiv_ours.dat};

\addlegendentry{\textsc{Dream-RSI}}

% ============================================================
% Recursive Fixed Exploration
% ============================================================
\addplot[
  const plot,
  color=baselinegray,
  line width=1.5pt,
  densely dashed,
  mark=square*,
  mark size=1.0pt,
  mark options={
    fill=white,
    draw=baselinegray,
    line width=0.65pt
  }
]
table[x=budget,y=score]
{Main_Text/raw_data/data/convdiv_baseline.dat};

\addlegendentry{Fixed Exploration}

% ============================================================
% Dream-RSI + History as Guidance
% ============================================================
\addplot[
  const plot,
  color=guidanceorange,
  line width=1.8pt,
  densely dashdotted,
  mark=triangle*,
  mark size=1.15pt,
  mark options={
    fill=guidanceorange,
    draw=guidanceorange
  }
]
table[x=budget,y=score]
{Main_Text/raw_data/data/convdiv_ours_guidance.dat};

\addlegendentry{\textsc{Dream-RSI} + Guidance}

% ============================================================
% Fixed Exploration + History as Guidance
% ============================================================
\addplot[
  const plot,
  color=guidancegreen,
  line width=1.6pt,
  densely dotted,
  mark=diamond*,
  mark size=1.05pt,
  mark options={
    fill=white,
    draw=guidancegreen,
    line width=0.65pt
  }
]
table[x=budget,y=score]
{Main_Text/raw_data/data/convdiv_baseline_guidance.dat};

\addlegendentry{Fixed + Guidance}

\end{axis}
\end{tikzpicture}

\vspace{-0.6em}
\caption{
Discovery performance on ConvDiv.
Using history as an interactive replay simulator outperforms using it only as guidance.
}
\label{fig:convdiv_guidance}
\vspace{-0.8em}
\end{wrapfigure}
We further investigate how the nature of the historical inductive bias affects long-horizon discovery. A natural alternative for utilizing history is to abstract prior trajectories into high-level directional insights, which are directly injected into the prompt as explicit semantic guidance for subsequent rounds. To evaluate the efficacy of this prompt-level semantic guidance, we apply it to both Recursive Fixed Exploration and \textsc{Dream-RSI}. As illustrated in Figure~\ref{fig:convdiv_guidance}, explicit directional guidance consistently underperforms its unguided counterpart across both paradigms under equivalent discovery budgets. These results suggest that in long-horizon discovery---where multiple parallel threads are deployed for exploration---imposing strong semantic inductive biases regarding future search directions tends to over-constrain the search space and impede diverse exploration.


\subsection{Analysis of Evolution of Exploration Behavior}
% ============================================================
% Preamble
% ============================================================
\pgfplotsset{compat=1.18}

\definecolor{oursblue}{HTML}{2563EB}
\definecolor{barblue}{HTML}{BFDBFE}

% ============================================================
% Data
% ============================================================

\pgfplotstableread{
round perf
0 0.427
1 0.625
2 0.855
3 1.403
4 1.488
5 1.499
6 1.770
7 1.880
8 1.898
}\perfdata

\pgfplotstableread{
round attempts
0 110
1 110
2 87
3 80
4 50
5 92
6 80
7 91
8 86
}\attemptdata

% ============================================================
% Figure
% ============================================================

\begin{wrapfigure}[18]{r}{0.5\linewidth}
\centering
\vspace{-2em}
% ------------------------------------------------------------
% (a) Round-best performance
% ------------------------------------------------------------
\begin{tikzpicture}
\begin{axis}[
    width=\linewidth,
    height=3.6cm,
    xmin=-0.4,
    xmax=8.4,
    ymin=0.35,
    ymax=2.03,
    xtick={0,1,2,3,4,5,6,7,8},
    xticklabels={E0,E1,E2,E3,E4,E5,E6,E7,E8},
    ylabel={Performance (1/ms)},
    axis lines=left,
    axis line style={black!55,line width=0.6pt},
    ymajorgrids=true,
    grid style={black!8,line width=0.3pt},
    tick align=outside,
    tick style={black!55},
    tick label style={font=\scriptsize},
    label style={font=\scriptsize},
    clip=false,
]

\addplot[
    color=oursblue,
    line width=1.5pt,
    mark=*,
    mark size=2.4pt,
    mark options={
        fill=white,
        draw=oursblue,
        line width=1.0pt
    }
]
table[x=round,y=perf] {\perfdata};

\addplot[
    only marks,
    mark=none,
    point meta=explicit symbolic,
    nodes near coords,
    every node near coord/.append style={
        font=\tiny\bfseries,
        text=oursblue,
        yshift=5pt
    }
]
table[
    x=round,
    y=perf,
    meta=perf
] {\perfdata};

\end{axis}
\end{tikzpicture}

\vspace{-2mm}

{\scriptsize\textbf{(a) Round-best performance}}

\vspace{1mm}

% ------------------------------------------------------------
% (b) Evaluated attempts
% ------------------------------------------------------------
\begin{tikzpicture}
\begin{axis}[
    width=\linewidth,
    height=3.6cm,
    xmin=-0.4,
    xmax=8.4,
    ymin=0,
    ymax=125,
    xtick={0,1,2,3,4,5,6,7,8},
    xticklabels={E0,E1,E2,E3,E4,E5,E6,E7,E8},
    ylabel={Evaluated attempts},
    xlabel={Recursive execution round},
    axis lines=left,
    axis line style={black!55,line width=0.6pt},
    ymajorgrids=true,
    grid style={black!8,line width=0.3pt},
    tick align=outside,
    tick style={black!55},
    tick label style={font=\scriptsize},
    label style={font=\scriptsize},
    clip=false,
]

% bars
\addplot[
    ybar,
    bar width=14pt,
    fill=barblue,
    draw=oursblue,
    line width=0.7pt
]
table[x=round,y=attempts] {\attemptdata};

% line connecting bar tops
\addplot[
    color=oursblue,
    line width=1.3pt,
    mark=*,
    mark size=2.2pt,
    mark options={
        fill=white,
        draw=oursblue,
        line width=1.0pt
    }
]
table[x=round,y=attempts] {\attemptdata};

% value labels
\addplot[
    only marks,
    mark=none,
    point meta=explicit symbolic,
    nodes near coords,
    every node near coord/.append style={
        font=\tiny\bfseries,
        text=oursblue,
        yshift=7pt
    }
]
table[
    x=round,
    y=attempts,
    meta=attempts
] {\attemptdata};

\end{axis}
\end{tikzpicture}

\vspace{-2mm}

{\scriptsize\textbf{(b) Exploration effort}}

\caption{
\textbf{Evolution of exploration behavior on ConvDiv.}
\textbf{(a)} Round-best performance across recursive execution rounds.
\textbf{(b)} The number of evaluated attempts in each round.
}
\label{fig:convdiv-dynamics}

\end{wrapfigure}
Figure~\ref{fig:convdiv-dynamics} illustrates how the learned exploration policy evolves across recursive rounds on ConvDiv. As shown, the exploration policy exhibits a clear adaptive pattern: as performance improves, it initially conserves discovery compute (e.g., reducing the number of evaluated attempts from 110 to 50). When progress subsequently plateaus, it increases exploration effort again, coinciding with further performance gains..